\documentclass[11pt,a4paper]{article}
\usepackage[T1]{fontenc}
\usepackage{lmodern}
\usepackage[margin=2.5cm]{geometry}
\usepackage{amsmath}
\usepackage{listings}
\usepackage{xcolor}
\usepackage{booktabs}
\usepackage{enumitem}
\usepackage{hyperref}

\lstset{
  language=C,
  basicstyle=\ttfamily\footnotesize,
  keywordstyle=\bfseries,
  commentstyle=\itshape\color{black!60},
  numbers=left,
  numberstyle=\tiny,
  numbersep=8pt,
  frame=single,
  breaklines=true,
  breakatwhitespace=false,
  postbreak=\mbox{$\hookrightarrow$\space},
  showstringspaces=false,
  columns=fullflexible,
  keepspaces=true,
  xleftmargin=1.5em,
  framexleftmargin=1.5em,
  morekeywords={bool,true,false}
}

% \code{file}{first}{last}: show real lines first..last of a source file, numbered as in the file
\newcommand{\code}[3]{\lstinputlisting[firstline=#2,lastline=#3,firstnumber=#2]{../#1}}

\setlist{itemsep=2pt}

\title{Mini Golf: A Line-by-Line Guide to the Code\\[4pt]
\large \texttt{types.h}, \texttt{input}, \texttt{level}, \texttt{putter} and the menu in \texttt{main.c}}
\author{Golf-Game project}
\date{\today}

\begin{document}
\maketitle
\tableofcontents
\newpage

%======================================================================
\section{Introduction}

This document walks through every line of five parts of the game:

\begin{itemize}
  \item \texttt{types.h} -- the data types the whole game shares;
  \item \texttt{input.h} / \texttt{input.c} -- reading the mouse;
  \item \texttt{level.h} / \texttt{level.c} -- loading maps, the course, scoring and the editor;
  \item \texttt{putter.h} / \texttt{putter\_simple.c} -- turning a mouse click into a shot;
  \item the menu part of \texttt{main.c}.
\end{itemize}

All code listings are taken directly from the source files, and the line numbers on the left are the
real line numbers in each file. So ``line 36 of \texttt{input.c}'' in the text is line 36 in your editor too.

The game uses the \textbf{raylib} library. A few raylib types appear everywhere:

\begin{center}
\begin{tabular}{@{}ll@{}}
\toprule
raylib type & meaning \\
\midrule
\texttt{Vector2}   & a 2D point or direction: \texttt{\{float x, float y\}} \\
\texttt{Rectangle} & \texttt{\{float x, float y, float width, float height\}}, $(x,y)$ is the top-left corner \\
\texttt{Camera2D}  & where the camera looks, its zoom and rotation \\
\texttt{Texture2D} & an image loaded onto the graphics card \\
\texttt{Music}     & a streamed audio file (mp3, ogg, \dots) \\
\bottomrule
\end{tabular}
\end{center}

The screen's $y$ axis points \emph{down}, so ``up'' on the screen is negative $y$.

%======================================================================
\section{\texttt{types.h}}

\texttt{types.h} has no functions. It only defines constants, enums and structs that every other
file uses, so all modules agree on what a ball, a wall or a hole looks like.

\subsection{Include guard and constants}

\code{types.h}{1}{10}

\begin{description}
  \item[Lines 1--2, 109.] This is an \emph{include guard}. The first time the compiler reads the file,
  \texttt{TYPES\_H} is not defined, so it defines it and reads the rest. If another header includes
  \texttt{types.h} again (for example \texttt{ball.h} and \texttt{level.h} both include it), \texttt{TYPES\_H}
  is already defined and the whole file is skipped up to the \texttt{\#endif} on line 109. Without this,
  every struct would be defined twice and the program would not compile.
  \item[Line 4.] Includes raylib so that \texttt{Vector2}, \texttt{Rectangle} and \texttt{Camera2D} are known.
  \item[Lines 7--10.] Maximum sizes. A hole can have at most 128 walls and 32 zones, and a course at most
  18 holes. These are used as array sizes, so no dynamic memory (\texttt{malloc}) is needed.
  \texttt{MAX\_OBSTACLES} is not used anywhere yet.
\end{description}

\subsection{\texttt{SurfaceType}}

\code{types.h}{12}{24}

An \texttt{enum} gives names to integers. The first name is 0 and each next one is one larger:

\begin{center}
\begin{tabular}{@{}clrl@{}}
\toprule
value & name & friction & behaviour (defined in \texttt{physics.c}) \\
\midrule
0 & \texttt{SURF\_GREEN}   & 400   & the default ground where there is no zone \\
1 & \texttt{SURF\_FAIRWAY} & 300   & a little faster than green \\
2 & \texttt{SURF\_SAND}    & 1400  & slows the ball a lot \\
3 & \texttt{SURF\_ICE}     & 90    & very slippery \\
4 & \texttt{SURF\_MUD}     & 2600  & strongest brake \\
5 & \texttt{SURF\_WATER}   & --    & hazard: the ball is dropped and gets $+1$ stroke \\
6 & \texttt{SURF\_BOOST}   & $-600$ & negative friction, so it speeds the ball up \\
\bottomrule
\end{tabular}
\end{center}

\begin{description}
  \item[Line 22.] \texttt{SURF\_COUNT} is a trick: because it comes last, its value is 7, which is exactly the
  number of real surfaces. Arrays like \texttt{kfriction[SURF\_COUNT]} and loops like
  \texttt{(zone\_type + 1) \% SURF\_COUNT} use it. If you add a new surface before it, the count updates itself.
  \item The numbers in the table are what you write after \texttt{zone} in a level file.
\end{description}

\subsection{\texttt{BallState} and \texttt{GameStateId}}

\code{types.h}{27}{46}

\begin{description}
  \item[Lines 28--35, \texttt{BallState}.] What the ball is doing right now:
  \begin{itemize}
    \item \texttt{BALL\_AIM} -- resting, the player may hit it;
    \item \texttt{BALL\_CHARGE} -- declared, but the putter does not use it (the putter has its own phase);
    \item \texttt{BALL\_ROLL} -- moving, physics runs every frame;
    \item \texttt{BALL\_SUNK} -- in the cup, the hole is finished;
    \item \texttt{BALL\_OOB} -- ``out of bounds'': in water or outside the field; it will be dropped.
  \end{itemize}
  \item[Lines 39--46, \texttt{GameStateId}.] The state of a round. \texttt{main.c} uses \texttt{GS\_PLAYING}
  (hitting the ball), \texttt{GS\_HOLE\_DONE} (the short pause showing ``BIRDIE'' etc.) and
  \texttt{GS\_SCOREBOARD} (after the last hole). \texttt{GS\_MENU} and \texttt{GS\_PAUSED} are not used: the
  menu is controlled by a separate \texttt{ScreenState} in \texttt{main.c} (Section~\ref{sec:menu}).
\end{description}

\subsection{\texttt{Ball}}

\code{types.h}{49}{58}

\begin{description}
  \item[Line 52.] \texttt{pos} -- the centre of the ball in world pixels.
  \item[Line 53.] \texttt{vel} -- the velocity in pixels per second. Each frame the position moves by
  \texttt{vel * dt}.
  \item[Line 54.] \texttt{radius} -- set to 7 pixels in \texttt{BallInit}.
  \item[Line 55.] \texttt{state} -- one of the \texttt{BallState} values above.
  \item[Line 56.] \texttt{strokes} -- shots taken on this hole, including penalty strokes.
  \item[Line 57.] \texttt{last\_safe\_pos} -- the last place the ball stopped normally. It is used as a fallback
  drop position if a hole has no drop zone.
\end{description}

\subsection{\texttt{Wall} and \texttt{Zone}}

\code{types.h}{61}{75}

\begin{description}
  \item[Line 64.] A wall is a solid rectangle.
  \item[Line 65.] \texttt{bounce} is the \emph{restitution}: how much speed the ball keeps when it hits the wall.
  0.72 keeps 72\,\%, 0.95 is almost a perfect bumper, 0 would stop the ball dead.
  \item[Line 72.] A zone is a rectangle of ground \dots
  \item[Line 73.] \dots with a surface type (sand, ice, \dots) \dots
  \item[Line 74.] \dots and an optional wind. \texttt{wind} is an \emph{acceleration} in pixels per second$^2$.
  \texttt{(0,-140)} pushes the ball upward a little more every frame while it is inside the zone.
\end{description}

\subsection{\texttt{Hole}}

\code{types.h}{78}{95}

\begin{description}
  \item[Lines 81--82.] A fixed array of 128 walls plus a counter. Only the first \texttt{wall\_count} entries are
  real walls, the rest are unused. This ``array + count'' pattern is how C stores a list whose length changes,
  without \texttt{malloc}.
  \item[Lines 84--85.] The same pattern for zones.
  \item[Line 87.] \texttt{number} -- the hole number shown in the HUD.
  \item[Line 88.] \texttt{par} -- how many strokes a good player needs.
  \item[Lines 89--91.] Where the ball starts, where the cup is, and the cup's radius.
  \item[Line 93.] \texttt{bounds} -- the whole playable field. If the ball's centre leaves it, the ball is out of bounds.
  \item[Line 94.] \texttt{drop\_zone} -- after water or out of bounds, the ball is placed at the centre of this rectangle.
\end{description}

\subsection{\texttt{Game}}

\code{types.h}{98}{109}

This struct was planned to hold the whole game, but \texttt{main.c} does not use it. The real game state is the
\texttt{GameApp} struct in \texttt{main.c} (Section~\ref{sec:gameapp}). Changing \texttt{Game} has no effect.

%======================================================================
\section{\texttt{input.h} and \texttt{input.c}}

The input module is the only place where the mouse is read during play. Every frame it produces one
\texttt{InputState} that the rest of the game uses.

\subsection{\texttt{input.h}}

\code{input.h}{1}{41}

\begin{description}
  \item[Lines 1--2, 41.] Include guard, as in \texttt{types.h}.
  \item[Lines 9--23, \texttt{InputState}.] Everything that happened with the mouse \emph{this frame}:
  \begin{itemize}
    \item \texttt{aim\_direction}, \texttt{aim\_angle} -- which way the player is aiming, as a unit vector and as an
    angle in radians;
    \item \texttt{pointer\_world} -- the mouse position in world coordinates;
    \item \texttt{pointer\_valid} -- \texttt{true} if the mouse is far enough from the ball to give a clean aim;
    \item \texttt{charge\_pressed}, \texttt{charge\_held}, \texttt{charge\_released} -- the left button went down this
    frame / is being held / came up this frame;
    \item \texttt{cancel} -- a quick right-click happened (cancels a shot);
    \item \texttt{confirm} -- a left click, used to skip the ``hole done'' pause;
    \item \texttt{cam\_zoom\_delta} -- meant for mouse-wheel zoom (see the note in Section~\ref{sec:inputpoll}).
  \end{itemize}
  \item[Lines 25--31, \texttt{InputSystem}.] Values that must be \emph{remembered between frames}: the last good
  aim angle and the timer and distance for the right button. One \texttt{InputSystem} lives inside \texttt{GameApp}
  for the whole game, while a new \texttt{InputState} is made every frame.
  \item[Lines 33--37.] The three functions the module offers.
\end{description}

\subsection{Constants, \texttt{DirFromAngle} and \texttt{InputInit}}

\code{input.c}{1}{25}

\begin{description}
  \item[Lines 1--4.] \texttt{raymath.h} gives vector helpers like \texttt{Vector2Subtract}; \texttt{math.h} gives
  \texttt{cosf}, \texttt{sinf}, \texttt{atan2f}.
  \item[Line 7.] If the mouse is closer than 40 pixels to the ball, the aim is not updated (explained below).
  \item[Lines 9--10.] A right-click counts as ``cancel'' only if it lasts less than 0.25 seconds and the mouse
  moves less than 6 pixels.
  \item[Line 12.] One notch of the mouse wheel would change the zoom by 0.1.
  \item[Lines 16--18, \texttt{DirFromAngle}.] Turns an angle $\theta$ into the unit vector $(\cos\theta, \sin\theta)$.
  Angle 0 points right, $\pi/2$ points down (because screen $y$ grows downward).
  \item[Lines 20--25, \texttt{InputInit}.] Puts every remembered value to zero. The \texttt{->} operator
  accesses a field through a pointer: \texttt{sys->aim\_angle} is the same as \texttt{(*sys).aim\_angle}.
\end{description}

\subsection{\texttt{InputPoll}}\label{sec:inputpoll}

\code{input.c}{27}{70}

\begin{description}
  \item[Line 27.] The function receives the remembered state \texttt{sys}, the ball's position, the camera and the
  frame time \texttt{dt}, and \emph{returns} a fresh \texttt{InputState} by value.
  \item[Line 28.] \texttt{= \{0\}} sets every field to zero / \texttt{false}. So anything not set below is ``nothing
  happened''.
  \item[Line 30.] \texttt{GetMousePosition()} gives the mouse in window pixels. Because the camera can move and
  zoom, that is not the same as world coordinates. \texttt{GetScreenToWorld2D} converts it using the camera.
  \item[Line 33.] The vector from the ball to the mouse.
  \item[Line 35.] If that vector is longer than 40 pixels \dots\ Comparing the \emph{squared} length with
  $40^2 = 1600$ gives the same answer as comparing the length with 40, but avoids a slow square root.
  \item[Line 36.] \dots\ the aim angle is updated. \texttt{atan2f(y, x)} returns the angle of the vector
  $(x,y)$ in radians, in the range $(-\pi, \pi]$, and it handles every quadrant correctly.
  \item[Lines 37--41.] Remember whether the pointer gave a valid aim. If the mouse is very close to the ball,
  tiny movements would swing the angle wildly, so the old angle is kept instead.
  \item[Lines 42--43.] Copy the (new or old) angle and its direction vector into the result.
  \item[Lines 45--47.] raylib offers three different questions about a button, and the putter needs all three:
  \texttt{IsMouseButtonPressed} is true only in the one frame the button goes down;
  \texttt{IsMouseButtonDown} is true in every frame while it is held;
  \texttt{IsMouseButtonReleased} is true only in the frame it comes up.
  (\texttt{MOUSE\_LEFT\_BUTTON} and \texttt{MOUSE\_BUTTON\_LEFT} are two names for the same button.)
  \item[Line 49.] \texttt{confirm} is also a left click, used when a hole is finished.
  \item[Lines 51--55.] When the right button goes down: reset the timer and distance and remember that it is down.
  \item[Lines 56--59.] While it stays down: add this frame's time to the timer, and store how far the mouse moved.
  \item[Lines 61--64.] When it comes up: it was a ``cancel'' only if it was short \emph{and} the mouse stayed still.
  This way a right-\emph{drag} (for example to move a camera later) does not cancel a shot.
  \item[Line 65.] Stores the mouse wheel movement times 0.1.
  \item[Line 68.] Return the filled \texttt{InputState}.
\end{description}

\paragraph{Two small bugs.}
\begin{itemize}
  \item \texttt{cam\_zoom\_delta} is declared \texttt{bool} (line 21 of \texttt{input.h}) but receives a
  \texttt{float} on line 65. Any non-zero value becomes \texttt{true} (1), so the real amount is lost. It should be
  declared \texttt{float}. Nothing uses it yet, so the game is not affected.
  \item On line 58, \texttt{=} stores only the \emph{last frame's} mouse movement. To measure the whole drag it
  should be \texttt{+=}.
\end{itemize}

%======================================================================
\section{\texttt{level.h} and \texttt{level.c}}\label{sec:level}

\subsection{\texttt{level.h}}

\code{level.h}{1}{67}

\begin{description}
  \item[Lines 8--12, \texttt{Course}.] All holes of the game, how many were loaded, and which one is being played
  (\texttt{current} starts at 0).
  \item[Lines 15--21, \texttt{Edit\_Tool}.] The five tools of the map editor: draw a wall, draw a zone, place the tee,
  place the cup, draw the drop zone.
  \item[Lines 25--32, \texttt{Edtior}.] The editor's state (the name has a typo, ``Edtior'', but it is used
  consistently so it works): current tool, which surface a new zone gets, whether the mouse is dragging and where
  the drag started, the snapping grid size, and which wall is under the mouse.
  \item[Lines 38--62.] The public functions. \texttt{save\_best\_scores}, \texttt{load\_best\_scores},
  \texttt{surface\_from\_name} and \texttt{surface\_name} are declared but never written; that is fine as long as
  nobody calls them.
\end{description}

\subsection{The level file format}

Each hole is a text file \texttt{levels/holeNN.txt}. One command per line; \texttt{\#} starts a comment.

\begin{center}
\begin{tabular}{@{}ll@{}}
\toprule
command & meaning \\
\midrule
\texttt{hole 7}                      & hole number \\
\texttt{par 4}                       & par \\
\texttt{bounds x y w h}              & the playable rectangle \\
\texttt{tee x y}                     & start position \\
\texttt{cup x y r}                   & cup centre and radius \\
\texttt{drop x y w h}                & drop zone \\
\texttt{wall x y w h [bounce]}       & a wall; \texttt{bounce} is optional (default 0.72) \\
\texttt{zone t x y w h [wx wy]}      & a zone of surface number \texttt{t}; wind is optional \\
\texttt{end}                         & ignored \\
\bottomrule
\end{tabular}
\end{center}

Example from \texttt{hole07.txt}:
\begin{lstlisting}[numbers=none]
tee 160 740
cup 1420 200 16
wall 420 24 40 520 0.8            # a canyon wall
zone 3 460 400 420 240 0 -140     # ice with wind blowing up
\end{lstlisting}

\subsection{\texttt{AddZone}, \texttt{Validate\_hole}, \texttt{AddWall}}

\code{level.c}{1}{42}

\begin{description}
  \item[Line 14--17, \texttt{AddZone}.] If the zone array is full, do nothing. Otherwise write the new zone into
  position \texttt{zone\_count} and \emph{then} increase \texttt{zone\_count} (that is what the post-increment
  \texttt{zone\_count++} does). The \texttt{(Zone)\{\dots\}} syntax is a \emph{compound literal}: it builds a whole
  struct in one expression.
  \item[Lines 21--31, \texttt{Validate\_hole}.] Checks a hole after loading. Two values are \emph{repaired}: a par
  outside 1--9 becomes 3, and a cup radius under 4 becomes 14. Four problems \emph{reject} the hole: the field is
  smaller than $100\times100$, the tee or the cup is outside the field (\texttt{CheckCollisionPointRec} tests whether a
  point is inside a rectangle), or there are no walls at all. The \texttt{path} parameter is not used.
  \item[Lines 34--42, \texttt{AddWall}.] Same idea as \texttt{AddZone}. Note the double brace \texttt{\{\{} on line 34
  and the extra \texttt{\}} on line 42: the body is wrapped in an extra, useless block. It is legal C and changes
  nothing, but one pair of braces can be removed.
\end{description}

\subsection{\texttt{Parseline}}

\code{level.c}{44}{86}

This function reads one line of a level file.

\begin{description}
  \item[Line 47.] Skip leading spaces by moving the \texttt{line} pointer forward.
  \item[Line 48.] Ignore comment lines and empty lines.
  \item[Lines 51--52.] Read the first word into \texttt{k}. \texttt{\%31s} reads at most 31 characters, so the
  32-byte buffer (31 letters + the terminating \texttt{\textbackslash 0}) can never overflow. \texttt{sscanf}
  returns how many values it read; if not 1, the line is ignored.
  \item[Lines 54--85.] Compare the word with each command using \texttt{strcmp} (which returns 0 when two strings
  are equal) and read the numbers that follow.
  \item[The \texttt{\%*s} trick.] Each \texttt{sscanf} reads the \emph{whole line} again. \texttt{\%*s} means ``read
  a word and throw it away'', so it skips the command name and the numbers go into the right variables.
  \item[Lines 73--75.] Optional values: \texttt{b} is set to 0.72 \emph{before} \texttt{sscanf}. If the line has only
  four numbers, \texttt{sscanf} stops early and \texttt{b} keeps 0.72.
  \item[Lines 78--81.] The same for zones: wind defaults to $(0,0)$. The zone number is read as an \texttt{int} and
  passed where a \texttt{SurfaceType} is expected; C converts an \texttt{int} to an enum automatically.
  \item[Lines 83--85.] \texttt{end} is accepted and does nothing. Any unknown word falls through all the
  \texttt{if}s and is silently ignored.
\end{description}

\subsection{\texttt{LoadHoleFromFile}}

\code{level.c}{88}{113}

\begin{description}
  \item[Line 89.] If the file does not exist, fail.
  \item[Line 91.] \texttt{LoadFileText} (raylib) reads the whole file into one string in memory.
  \item[Line 94.] \texttt{*h = (Hole)\{0\};} clears every field of the hole, so no walls or zones from a previous map
  are left behind.
  \item[Lines 97--99.] Default values in case the file does not set them.
  \item[Line 102.] \texttt{strtok} cuts a string into pieces. The first call gets the text and the separators
  (\texttt{\textbackslash r} and \texttt{\textbackslash n}, so both Windows and Mac/Linux line endings work) and
  returns the first line.
  \item[Lines 104--108.] Each later call with \texttt{NULL} continues where the last one stopped and returns the next
  line, until it returns \texttt{NULL} at the end. Each line goes to \texttt{Parseline}.
  \item[Line 109.] Free the memory that \texttt{LoadFileText} allocated.
  \item[Line 111.] Return whether the hole is valid.
\end{description}

\subsection{\texttt{SaveHoleToFile}}

\code{level.c}{117}{124}

Not written yet. It always returns \texttt{false}, and \texttt{buf} and \texttt{n} are unused. It is meant to write a
hole back to a text file for the editor.

\subsection{The course}

\code{level.c}{127}{150}

\begin{description}
  \item[Lines 127--139, \texttt{course\_load}.] Tries \texttt{levels/hole01.txt}, \texttt{hole02.txt}, \dots
  \texttt{TextFormat} works like \texttt{sprintf}; \texttt{\%02d} prints the number with at least two digits and a
  leading zero (1 $\to$ \texttt{01}).
  \begin{itemize}
    \item Line 133: the loop \emph{stops} at the first missing file, so hole numbers must not have gaps.
    \item Line 135: a file that fails validation is skipped, but loading continues with the next number.
    Valid holes are stored one after another, so \texttt{holes[0]}, \texttt{holes[1]}, \dots\ are always filled.
    \item The loop runs for \texttt{i = 1 \dots\ 17}, so at most 17 files are tried.
  \end{itemize}
  \item[Lines 141--145, \texttt{course\_advance}.] Go to the next hole. Returns \texttt{false} if the current hole was the
  last one, which tells \texttt{main.c} to show the scoreboard.
  \item[Lines 147--150, \texttt{course\_current}.] Returns a \emph{pointer} to the current hole, so the caller works
  on the real hole and not a copy. If \texttt{current} is somehow out of range, it safely returns hole 0.
\end{description}

\subsection{Scoring}

\code{level.c}{153}{189}

\begin{description}
  \item[Lines 155--168, \texttt{score\_name}.] One stroke is always ``HOLE IN ONE!!!''. Otherwise
  \texttt{strokes - par} picks the name: $-3$ albatross, $-2$ eagle, $-1$ birdie, 0 par, $+1$ bogey, $+2$ double bogey,
  $+3$ triple bogey. The \texttt{default} case handles everything else: better than $-3$ gives ``GREAT'', worse than
  $+3$ gives the insult. (\texttt{case - 3} with a space is simply the number $-3$.)
  \item[Lines 171--179, \texttt{Course\_total}.] Adds up all scores. A score of 0 means ``not played yet'' and is
  skipped.
  \item[Lines 182--189, \texttt{course\_to\_par}.] Adds up \texttt{score - par} for played holes. A result of $-2$ means
  two under par for the round so far.
\end{description}

\subsection{The editor}

\code{level.c}{195}{246}

\begin{description}
  \item[Lines 197--203, \texttt{EditorInit}.] Start with the wall tool, sand as the zone surface, a 10-pixel grid and
  nothing selected ($-1$ means ``no wall'').
  \item[Lines 205--208, \texttt{Snap2Grid}.] Rounds a point to the nearest grid line. With a grid of 10,
  $x = 347$ becomes $\operatorname{round}(34.7)\cdot 10 = 350$. A grid of 1 means no snapping.
  \item[Lines 210--214, \texttt{RectFromCorners}.] The user may drag in any direction. \texttt{fminf} takes the
  smaller coordinates as the top-left corner and \texttt{fabsf} (absolute value) gives a positive width and height.
  \item[Lines 217--223, \texttt{RemoveWall}.] Deletes wall number \texttt{index} by shifting every later wall one place
  to the left, then reducing the count. Not called anywhere yet.
  \item[Lines 225--230, \texttt{WallAtPoint}.] Finds the wall under a point. It searches from the \emph{last} wall
  backwards, so if walls overlap, the one drawn on top is found first. Returns $-1$ if there is none.
  \item[Lines 233--234.] A global copy of one hole, used for a one-step undo.
  \item[Lines 236--239, \texttt{PushUndo}.] Saves a copy of the hole before a change.
  \item[Lines 241--246, \texttt{PopUndo}.] Swaps the current hole with the saved copy. Because it swaps (using
  \texttt{temporary}), calling it twice is a redo. Not called anywhere yet.
\end{description}

\code{level.c}{248}{313}

\begin{description}
  \item[Line 249.] The mouse position snapped to the grid.
  \item[Line 250.] Remember which wall is under the mouse (for highlighting).
  \item[Lines 254--258.] Keys 1--5 choose the tool.
  \item[Lines 262--264.] \texttt{Tab} cycles the zone surface: $0 \to 1 \to \dots \to 6 \to 0$. The \texttt{\%}
  (remainder) operator makes it wrap around, and \texttt{SURF\_COUNT} is why it wraps at the right number.
  \item[Lines 267--272.] \texttt{G} cycles the grid: $1 \to 5 \to 10 \to 20 \to 1$.
  \item[Lines 275--279.] When the left button goes down: save an undo copy and remember where the drag started.
  \item[Lines 281--283.] When it comes up: build the rectangle between the start and the current point.
  \item[Lines 285--307.] Use it according to the tool. Tee and cup use only the release point. Walls, zones and the
  drop zone need a rectangle at least $4\times4$ pixels, so a plain click does not create an invisible object.
\end{description}

\paragraph{Status.} \texttt{main.c} calls \texttt{EditorInit} but never calls \texttt{editor\_update}, so the
editor is not active in the game yet, even though the HUD mentions ``F1 editor''.

%======================================================================
\section{\texttt{putter.h} and \texttt{putter\_simple.c}}

\subsection{\texttt{putter.h}}

\code{putter.h}{1}{41}

\begin{description}
  \item[Lines 16--22, \texttt{PutterPhase}.] The phases of a swing. The simple putter uses only
  \texttt{putter\_idle} (waiting) and \texttt{putter\_backswing} (charging power). \texttt{strike},
  \texttt{follow} and \texttt{recover} were for the animated club in \texttt{unused/putter\_anim.c}.
  \item[Lines 24--33, \texttt{Putter}.]
  \begin{itemize}
    \item \texttt{phase} -- the current phase;
    \item \texttt{t}, \texttt{offset} -- only used by the animated version;
    \item \texttt{power} -- the power bar, from 0 to 1;
    \item \texttt{launch\_power} -- the power actually used for the last shot;
    \item \texttt{aim\_angle} -- a copy of the input aim;
    \item \texttt{dir} -- $+1$ while the bar fills, $-1$ while it empties;
    \item \texttt{fired} -- \texttt{true} in the frame a shot was taken.
  \end{itemize}
  \item[Lines 35--39.] The functions. Parameters marked \texttt{const} are only read, never changed.
\end{description}

\subsection{\texttt{putter\_simple.c}}

\code{putter_simple.c}{1}{24}

\begin{description}
  \item[Line 7.] Filling the bar from 0 to 1 takes 1.15 seconds.
  \item[Line 8.] The weakest possible shot is 5\,\% power, so a very quick click still moves the ball.
  \item[Lines 11--16, \texttt{putter\_init}.] Start idle with everything at zero. Line 13 is a chained assignment:
  it is read from right to left, so all five fields become 0.
  \item[Lines 18--20.] The putter is ``charging'' exactly when it is in the backswing. The HUD calls this to decide
  whether to draw the power bar.
  \item[Lines 22--24.] A simple getter for the current power.
\end{description}

\code{putter_simple.c}{26}{69}

\begin{description}
  \item[Line 27.] Always follow the player's aim.
  \item[Line 29.] The ball can only be hit when it is resting.
  \item[Lines 31--41, idle phase.] Keep the power at zero. If the ball is resting and the left button was
  \emph{just pressed}, switch to the backswing with the bar filling. \texttt{return} ends the function here, so the
  charging code below does not run in the same frame.
  \item[Lines 44--48, charging.] While the button is held, change the power by
  $\texttt{dir}\cdot dt / 1.15$. At 60 FPS, $dt \approx 0.0167$\,s, so the power grows by about 0.0145 per frame.
  When it passes 1 it is clamped to 1 and \texttt{dir} flips to $-1$; when it drops below 0 it flips back.
  So the bar goes up and down like a ping-pong ball until the player lets go.
  \item[Lines 50--54, cancel.] A quick right-click throws the charge away and returns to idle.
  \item[Lines 56--62, firing.] When the left button is released:
  \begin{itemize}
    \item line 57: use the current power, but at least \texttt{MIN\_POWER} (the \texttt{? :} is the conditional
    operator: \texttt{a ? b : c} means ``if \texttt{a} then \texttt{b} else \texttt{c}'');
    \item line 59: \texttt{BallLaunch} in \texttt{ball.c} sets the ball's velocity to
    \[ \vec v = (\cos\theta,\ \sin\theta)\cdot \texttt{power}\cdot 1300\ \text{px/s}, \]
    switches the ball to \texttt{BALL\_ROLL} and adds one stroke;
    \item lines 60--61: reset the bar and go back to idle.
  \end{itemize}
  \item[Lines 65--69, \texttt{putter\_draw}.] Draws nothing (there is no club picture). \texttt{(void)p;} tells the
  compiler ``I know this parameter is unused'', which removes a warning.
\end{description}

\paragraph{Example.} Holding for 0.575\,s gives power $0.575/1.15 = 0.5$, so the ball starts at 650\,px/s.
Holding for 1.6\,s: the bar reached 1 at 1.15\,s and has been falling for 0.45\,s, so power $= 1 - 0.45/1.15
\approx 0.61$.

\paragraph{How far does it roll?} With constant friction $a$, a ball starting at speed $v$ stops after
$d = v^2/(2a)$. A full shot on the fairway ($a=300$) rolls $1300^2/600 \approx 2800$\,px; on mud ($a=2600$) only
about 325\,px.

%======================================================================
\section{The menu in \texttt{main.c}}\label{sec:menu}

\subsection{\texttt{ScreenState}}

\code{main.c}{1}{22}

\begin{description}
  \item[Lines 1--10.] \texttt{main.c} includes every module because it connects them all.
  \item[Line 13.] A hole ends automatically after \texttt{par + 5} strokes.
  \item[Line 14.] The ``BIRDIE!'' message stays for 1.6 seconds.
  \item[Lines 17--22.] Which screen is on: the main menu, the game, the tutorial picture or the about picture. This is
  the top-level state; \texttt{GameStateId} (from \texttt{types.h}) only matters when the screen is
  \texttt{SCREEN\_GAME}.
\end{description}

\subsection{\texttt{GameApp}}\label{sec:gameapp}

\code{main.c}{25}{47}

\begin{description}
  \item[Lines 26--36.] Game data: round state, the course, the ball, the putter, the input memory, the editor, the
  camera, the score of each hole, the best score of each hole, the timer for the ``hole done'' pause, and
  \texttt{prev\_speed} (not used yet).
  \item[Lines 39--46.] Menu data: the current screen, three pictures for the menu, tutorial and about screens, the
  wind animation, the two music tracks, and whether sound is muted.
\end{description}

\subsection{\texttt{DrawFullscreenTexture}}

\code{main.c}{50}{56}

\texttt{DrawTexturePro} copies the \texttt{source} rectangle of the picture (here the whole picture) into the
\texttt{destination} rectangle on screen (here the whole window). If the sizes differ, the image is stretched, so
the background always fills the window, even after resizing. \texttt{origin} is the rotation point and
\texttt{0.0f} is the rotation angle. \texttt{WHITE} as the tint means ``draw the original colours''.

\subsection{\texttt{DrawMenuButton}}

\code{main.c}{59}{76}

\begin{description}
  \item[Lines 61--62.] Is the mouse over the button? (Menu buttons use screen coordinates, so no camera conversion is
  needed.)
  \item[Lines 64--66.] A lighter green when hovered, a darker green otherwise. The fourth number in a \texttt{Color} is
  alpha (opacity, 0--255).
  \item[Lines 68--69.] Draw the rounded button and its light outline. \texttt{0.20f} is the roundness and \texttt{10}
  the number of segments used to draw each corner.
  \item[Lines 71--73.] Centre the text: \texttt{MeasureText} gives the text width in pixels, so the text starts at
  $\text{centre} - \text{width}/2$. Vertically it starts half a font size above the centre.
  \item[Lines 74--75.] Return \texttt{true} only in the frame the button is clicked.
\end{description}

This style is called an \emph{immediate-mode} GUI: there are no button objects stored anywhere. Each frame the
function is called again, draws the button and answers ``was I clicked right now?''. To add a button you just add one
more call.

\subsection{\texttt{StartHole}}

\code{main.c}{79}{86}

Used when a new hole begins: put the ball on the tee (\texttt{BallInit} also resets the stroke count to 0), reset the
putter, restart the hole timer, jump the camera to the ball, and go back to \texttt{GS\_PLAYING}.

\subsection{Music functions}

\code{main.c}{89}{115}

\begin{description}
  \item[Lines 89--93, \texttt{StartGameMusic}.] Stop the menu song and play the game song, unless muted.
  \item[Lines 96--100, \texttt{StartMenuMusic}.] The opposite.
  \item[Lines 103--115, \texttt{ToggleMute}.] Flip \texttt{muted} with \texttt{!}, then set both volumes to 0 or back to
  normal. The music keeps playing silently, so un-muting continues the song instead of restarting it.
\end{description}

\subsection{\texttt{GameInit}}

\code{main.c}{118}{147}

\begin{description}
  \item[Line 120.] Clear the whole struct.
  \item[Line 122.] Load all holes from the \texttt{levels} folder.
  \item[Lines 126--131.] Initialise each module and set the round state to playing.
  \item[Line 132.] The game starts on the \textbf{main menu}, not in play.
  \item[Lines 135--141.] Load the pictures, the wind animation and the two songs. Line 139 makes the wind sprite use smooth (bilinear) scaling. The paths are relative to the folder the program is
  \emph{started from}, so run it from the project folder (\texttt{make run} or \texttt{./golf}). A file must also really
  be the format its name says: a JPEG renamed to \texttt{.png} will not load.
  \item[Lines 142--145.] Both songs loop, with the menu slightly louder.
  \item[Line 146.] Start the menu music.
\end{description}

\subsection{\texttt{DrawMainMenu}}

\code{main.c}{150}{178}

\begin{description}
  \item[Line 152.] The background picture fills the window.
  \item[Line 153.] A black rectangle with alpha 70 (about 27\,\% opaque) over the whole window darkens the picture so
  the white text is easier to read.
  \item[Lines 154--157.] The title ``MINI GOLF'' in size 70, centred horizontally, 80 pixels from the top.
  \item[Lines 158--165.] Five buttons, $260\times60$ pixels, all centred horizontally, with their tops at
  $y = 220, 300, 380, 460, 540$ (80 pixels apart). \texttt{centerX} is recomputed every frame, so the buttons stay
  centred when the window is resized.
  \item[Lines 166--169, PLAY.] Switch to the game screen and the game music.
  \item[Lines 170--171.] TUTORIAL and ABOUT US switch screens.
  \item[Lines 172--175, SOUND.] The label is chosen from \texttt{muted} each frame, so it always shows the current
  setting. Clicking it calls \texttt{ToggleMute}.
  \item[Line 176, EXIT.] Returns \texttt{true}; \texttt{main} uses this to end the loop.
  \item[Line 177.] Otherwise return \texttt{false}: keep running.
\end{description}

\subsection{Tutorial and About screens}

\code{main.c}{181}{194}

Each screen shows one full-window picture and a BACK button in the top-left corner (at $x=30$, $y=30$,
$160\times55$ pixels). Clicking BACK returns to the main menu.

\subsection{The main loop: choosing the screen}

\code{main.c}{197}{246}

\begin{description}
  \item[Line 199.] Ask for 4$\times$ anti-aliasing (smoother edges) and a resizable window.
  \item[Lines 202--203.] Open the window and the audio device. Audio must be started before any music is loaded.
  \item[Line 204.] By default raylib closes the window when \texttt{Esc} is pressed. \texttt{SetExitKey(KEY\_NULL)}
  turns that off, so \texttt{Esc} can be used for ``back to menu''.
  \item[Line 205.] Run at 60 frames per second.
  \item[Lines 207--209.] Create and initialise the game. \texttt{shouldQuit} becomes \texttt{true} when EXIT is clicked.
  \item[Line 212.] Keep looping until the window's close button is pressed or EXIT is clicked.
  \item[Lines 213--214.] Frame time, clamped to 0.05\,s so that one slow frame cannot move the ball too far.
  \item[Lines 215--216.] raylib streams music in small pieces. \texttt{UpdateMusicStream} must be called every frame
  for each song, otherwise the music stutters and stops.
  \item[Lines 217--223, main menu.] Draw the menu. \texttt{continue} jumps straight to the next loop iteration, so none
  of the game code below runs while the menu is open. The game is frozen, and PLAY resumes it exactly where it was.
  \item[Lines 224--239.] The tutorial and about screens work the same way; \texttt{Esc} also goes back to the menu.
  \item[Lines 240--246.] From here on the screen is \texttt{SCREEN\_GAME}. The input is read, and \texttt{Esc} returns
  to the menu and switches back to the menu music.
\end{description}

The screens and how you move between them:

\begin{center}
\begin{tabular}{@{}lll@{}}
\toprule
from & action & to \\
\midrule
main menu & PLAY          & game \\
main menu & TUTORIAL      & tutorial \\
main menu & ABOUT US      & about \\
main menu & EXIT          & program ends \\
tutorial / about & BACK or \texttt{Esc} & main menu \\
game      & \texttt{Esc}  & main menu \\
\bottomrule
\end{tabular}
\end{center}

\subsection{Cleaning up}

\code{main.c}{309}{322}

After the loop ends, the music is stopped, every song and picture is unloaded from memory, and the audio device
and window are closed. Every \texttt{Load\dots} call has a matching \texttt{Unload\dots} call.

\end{document}
